\section{Confrontation with Gravitational-Wave Observations}
  \label{sec:gw_constraints}

  \subsection{Mapping to the LVK Parametrization}
    \label{subsec:mapping-to-the-lvk-parametrization}

    Gravitational-wave catalogues constrain modified dispersion relations
    of the general form
    \[
      \omega^2 = c^2 k^2 + A_\alpha k^\alpha,
    \]
    with $\alpha=4$ corresponding to the leading higher-derivative
    correction compatible with locality and parity.

    The spectral Lorentzian completion derived above predicts
    \[
      \omega^2 = c^2 k^2 - \gamma \ell_\chi^2 k^4 + \mathcal{O}(k^6),
    \]
    with
    \[
      \gamma = \frac{1}{180\,\zeta}.
    \]

    Identifying coefficients gives
    \[
      A_4 = -\gamma \ell_\chi^2.
    \]

    Current analyses place bounds on $|A_4|$ from phase corrections
    accumulated during propagation over cosmological distances.
    Denoting the observational upper bound by $A_4^{\mathrm obs}$,
    the spectral prediction implies
    \[
      \ell_\chi \lesssim \sqrt{\frac{A_4^{\mathrm obs}}{\gamma}} = \sqrt{180\,\zeta\,A_4^{\mathrm obs}}.
    \]

    In the natural spectral scheme $\zeta=1/6$,
    this reduces to
    \[
      \ell_\chi \lesssim \sqrt{30\,A_4^{\mathrm obs}}.
    \]

  \subsection{Robustness and Scheme Dependence}
    \label{subsec:robustness-and-scheme-dependence}

    The coefficient $\zeta=\mathcal{O}(1)$ captures normalization
    conventions in the identification of the Einstein term and possible
    multiplicity factors associated with the underlying operator.
    While its precise numerical value is scheme-dependent,
    the following features are robust:

    \begin{itemize}
      \item The sign of the $k^4$ correction is negative.
      \item The structure of the dispersion is fixed by the ratio of Seeley--DeWitt coefficients.
      \item The constraint on $\ell_\chi$ scales as $\sqrt{\zeta}$ and is therefore stable up to factors of order unity.
    \end{itemize}

    Thus, gravitational-wave observations provide a direct infrared probe
    of the pre-geometric scale introduced in the spectral framework.

  \subsection{Luminal Front Propagation}
    \label{subsec:luminal-front-propagation}

    Multi-messenger observations constrain the difference between the
    speed of gravitational waves and the speed of light to extremely high precision.
    In the present framework, the characteristic operator remains
    $\Box$ at leading order, and the correction enters only at order $k^4$.
    The front velocity is therefore luminal, and deviations of the group
    velocity are suppressed by $(k\ell_\chi)^2$.

    Consequently, the Lorentzian spectral completion is consistent with
    current bounds on gravitational-wave speed while predicting a specific
    higher-order dispersive correction.

  \subsection{Summary of Observational Implications}
    \label{subsec:summary-of-observational-implications}

    Modified dispersion relations for gravitational waves are
    constrained by the analyses of the \cite{LIGOScientificCollaboration2017}
    and subsequent catalogues of the \cite{LIGOScientificCollaboration,VirgoCollaboration}.
    Multi-messenger observations further constrain deviations from
    luminal propagation~\cite{LIGOScientificCollaboration2017GW170817}.

    The Lorentzian completion of the spectral entropy functional leads to:

    \begin{enumerate}
      \item Exactly two propagating tensor polarizations in the infrared.
      \item A modified dispersion relation with fixed structure $\omega^2 = c^2 k^2 - \gamma \ell_\chi^2 k^4$.
      \item A direct observational bound on the pre-geometric scale
      $\ell_\chi$ from gravitational-wave data.
    \end{enumerate}

    This establishes, for the first time within the spectral framework,
    a concrete and testable connection between the emergent geometric dynamics and interferometric observations.
